% --- Archon physics formalization source begin ---
% archon:physics
% archon:covers PhyXMiniProblems/problem_phyx_mini_0425.lean
% archon:source-report reports/phyx_mini/problem_phyx_mini_0425.source.json

\chapter{Physics problem phyx\_mini\_0425}
\label{ch:PhyXMiniProblems_problem_phyx_mini_0425}

\paragraph{Problem source.}
\#\# Physical scenario

A heat engine using \textbackslash{}( 1.0\textbackslash{},\textbackslash{}mathrm\{mol\} \textbackslash{}) of a monatomic gas follows the cycle shown in figure. \textbackslash{}( 3750\textbackslash{},\textbackslash{}mathrm\{J\} \textbackslash{}) of heat energy is transferred to the gas during process \textbackslash{}( 1 \textbackslash{}rightarrow 2 \textbackslash{}).

\#\# Question

What is the thermal efficiency of this heat engine?

\#\# Answer choices

- A: A: 0.46

- B: B: 0.40

- C: C: 0.67

- D: D: 0.15

\#\# Figure caption (auxiliary; use the image as primary evidence)

The image shows a pressure-volume (\textbackslash{}( p \textbackslash{})-\textbackslash{}( V \textbackslash{})) diagram depicting a thermodynamic cycle with a square loop. The axes are labeled with pressure \textbackslash{}( p \textbackslash{}) in kilopascals (kPa) on the vertical axis and volume \textbackslash{}( V \textbackslash{}) on the horizontal axis.

1. **Points and Cycle**: The cycle consists of four points labeled 1, 2, 3, and 4 connected by arrows indicating the direction of the cycle:
   - **1 to 2**: Vertical line with an upward arrow.
   - **2 to 3**: Horizontal line with a rightward arrow.
   - **3 to 4**: Vertical line with a downward arrow.
   - **4 to 1**: Horizontal line with a leftward arrow.

2. **Textual Information**:
   - Between points 1 and 2, there's a horizontal purple arrow labeled "3750 J".
   - Near point 1, the temperature is indicated as \textbackslash{}( T\_1 = 300 \textbackslash{}, K \textbackslash{}).

3. **Axes Labels**:
   - The pressure axis is marked with \textbackslash{}( p \textbackslash{}max \textbackslash{}) near the top left and increments marked at 300 kPa.
   - The volume axis is marked as \textbackslash{}( V\_\{\textbackslash{}min\} \textbackslash{}), \textbackslash{}( V\_\{\textbackslash{}max\} = 2V\_\{\textbackslash{}min\} \textbackslash{}).

The diagram visually represents a constant volume and constant pressure process within the thermodynamic cycle.

\#\# Dataset metadata

Laws of Thermodynamics; Physical Model Grounding Reasoning, Multi-Formula Reasoning

\paragraph{Recorded answer/context.}
D: D: 0.15

\paragraph{Figure/image path.}
/root/proposal\_for\_physic/science-mango/phyx\_mini\_run/phyx\_data/test\_image/425.png

\paragraph{Formalization target.}
create a compiling Lean file with sorry bodies at `PhyXMiniProblems/problem\_phyx\_mini\_0425.lean`. The Lean declarations must preserve the physical quantities, dimensions or dimensional roles, figure labels, governing-law hypotheses, and final relation expressed by this problem.
Use Mathlib/Physlib names found through LeanExplore where available. If a physics API is missing, introduce faithful local abstractions rather than scalar placeholder aliases.

\begin{theorem}[Physics formalization target]
\label{thm:physics:phyx_mini_0425:target}
\lean{PhyXMiniProblems.ProblemPhyXMini0425.thermalEfficiency_eq_answerChoiceD}
\uses{def:physics:phyx-mini-0425:phyxminiproblems-problemphyxmini0425-monatomicrectangularheatengine, def:physics:phyx-mini-0425:phyxminiproblems-problemphyxmini0425-matchesproblemstatement, def:physics:phyx-mini-0425:phyxminiproblems-problemphyxmini0425-matchesprimarypvdiagram, def:physics:phyx-mini-0425:phyxminiproblems-problemphyxmini0425-usesclassroomsigasconstant, def:physics:phyx-mini-0425:phyxminiproblems-problemphyxmini0425-hasphysicalcycleparameters, def:physics:phyx-mini-0425:phyxminiproblems-problemphyxmini0425-satisfiesmonatomicidealgaslaws, def:physics:phyx-mini-0425:phyxminiproblems-problemphyxmini0425-thermalefficiency, def:physics:phyx-mini-0425:phyxminiproblems-problemphyxmini0425-displayedefficiency, def:physics:phyx-mini-0425:phyxminiproblems-problemphyxmini0425-recordedanswerchoice, def:physics:phyx-mini-0425:phyxminiproblems-problemphyxmini0425-roundstodisplayedhundredth}
The assigned autoformalize agent should translate this physics problem into Lean declarations in the covered file, with theorem and lemma proof bodies written as `by sorry`.
\end{theorem}
\begin{proof}
This is an autoformalization task, not a proof task. Produce statements that can later be proved by the physics prover without weakening the physical model.
\end{proof}

% --- Archon declaration topology begin ---
\section{Declaration topology}
\begin{definition}[Volume Quantity]
  \label{def:physics:phyx-mini-0425:phyxminiproblems-problemphyxmini0425-volumequantity}
  \lean{PhyXMiniProblems.ProblemPhyXMini0425.VolumeQuantity}
  A nonnegative physical volume carrying dimension L³
\end{definition}
\begin{proof}
  This is the declared model definition of Volume Quantity.
\end{proof}
\begin{definition}[volume In Cubic Meters]
  \label{def:physics:phyx-mini-0425:phyxminiproblems-problemphyxmini0425-volumeincubicmeters}
  \lean{PhyXMiniProblems.ProblemPhyXMini0425.volumeInCubicMeters}
  \uses{def:physics:phyx-mini-0425:phyxminiproblems-problemphyxmini0425-volumequantity}
  Cubic-metre readout of a physical volume
\end{definition}
\begin{proof}
  This is the declared model definition of volume In Cubic Meters.
\end{proof}
\begin{definition}[pressure In Pascals]
  \label{def:physics:phyx-mini-0425:phyxminiproblems-problemphyxmini0425-pressureinpascals}
  \lean{PhyXMiniProblems.ProblemPhyXMini0425.pressureInPascals}
  Pascal readout of a physical pressure
\end{definition}
\begin{proof}
  This is the declared model definition of pressure In Pascals.
\end{proof}
\begin{definition}[pressure In Kilopascals]
  \label{def:physics:phyx-mini-0425:phyxminiproblems-problemphyxmini0425-pressureinkilopascals}
  \lean{PhyXMiniProblems.ProblemPhyXMini0425.pressureInKilopascals}
  \uses{def:physics:phyx-mini-0425:phyxminiproblems-problemphyxmini0425-pressureinpascals}
  Kilopascal readout of a physical pressure
\end{definition}
\begin{proof}
  This is the declared model definition of pressure In Kilopascals.
\end{proof}
\begin{definition}[energy In Joules]
  \label{def:physics:phyx-mini-0425:phyxminiproblems-problemphyxmini0425-energyinjoules}
  \lean{PhyXMiniProblems.ProblemPhyXMini0425.energyInJoules}
  Joule readout of a signed physical energy, heat transfer, or work
\end{definition}
\begin{proof}
  This is the declared model definition of energy In Joules.
\end{proof}
\begin{definition}[temperature In Kelvin]
  \label{def:physics:phyx-mini-0425:phyxminiproblems-problemphyxmini0425-temperatureinkelvin}
  \lean{PhyXMiniProblems.ProblemPhyXMini0425.temperatureInKelvin}
  Kelvin readout of a Physlib absolute temperature stored in an explicitly chosen zero-preserving temperature unit
\end{definition}
\begin{proof}
  This is the declared model definition of temperature In Kelvin.
\end{proof}
\begin{definition}[Cycle State]
  \label{def:physics:phyx-mini-0425:phyxminiproblems-problemphyxmini0425-cyclestate}
  \lean{PhyXMiniProblems.ProblemPhyXMini0425.CycleState}
  The four numbered vertices printed on the pressure--volume loop
\end{definition}
\begin{proof}
  This is the declared model definition of Cycle State.
\end{proof}
\begin{definition}[Cycle Leg]
  \label{def:physics:phyx-mini-0425:phyxminiproblems-problemphyxmini0425-cycleleg}
  \lean{PhyXMiniProblems.ProblemPhyXMini0425.CycleLeg}
  The four directed sides indicated by arrows in the supplied image
\end{definition}
\begin{proof}
  This is the declared model definition of Cycle Leg.
\end{proof}
\begin{definition}[leg Source]
  \label{def:physics:phyx-mini-0425:phyxminiproblems-problemphyxmini0425-legsource}
  \lean{PhyXMiniProblems.ProblemPhyXMini0425.legSource}
  \uses{def:physics:phyx-mini-0425:phyxminiproblems-problemphyxmini0425-cyclestate, def:physics:phyx-mini-0425:phyxminiproblems-problemphyxmini0425-cycleleg}
  Initial state of each directed leg
\end{definition}
\begin{proof}
  This is the declared model definition of leg Source.
\end{proof}
\begin{definition}[leg Target]
  \label{def:physics:phyx-mini-0425:phyxminiproblems-problemphyxmini0425-legtarget}
  \lean{PhyXMiniProblems.ProblemPhyXMini0425.legTarget}
  \uses{def:physics:phyx-mini-0425:phyxminiproblems-problemphyxmini0425-cyclestate, def:physics:phyx-mini-0425:phyxminiproblems-problemphyxmini0425-cycleleg}
  Final state of each directed leg
\end{definition}
\begin{proof}
  This is the declared model definition of leg Target.
\end{proof}
\begin{definition}[Process Kind]
  \label{def:physics:phyx-mini-0425:phyxminiproblems-problemphyxmini0425-processkind}
  \lean{PhyXMiniProblems.ProblemPhyXMini0425.ProcessKind}
  Thermodynamic constraint represented by a side of the rectangle
\end{definition}
\begin{proof}
  This is the declared model definition of Process Kind.
\end{proof}
\begin{definition}[Segment Orientation]
  \label{def:physics:phyx-mini-0425:phyxminiproblems-problemphyxmini0425-segmentorientation}
  \lean{PhyXMiniProblems.ProblemPhyXMini0425.SegmentOrientation}
  Geometric orientation of a straight segment in the p-V plane
\end{definition}
\begin{proof}
  This is the declared model definition of Segment Orientation.
\end{proof}
\begin{definition}[displayed Process Kind]
  \label{def:physics:phyx-mini-0425:phyxminiproblems-problemphyxmini0425-displayedprocesskind}
  \lean{PhyXMiniProblems.ProblemPhyXMini0425.displayedProcessKind}
  \uses{def:physics:phyx-mini-0425:phyxminiproblems-problemphyxmini0425-cycleleg, def:physics:phyx-mini-0425:phyxminiproblems-problemphyxmini0425-processkind}
  Process kinds read from the primary image
\end{definition}
\begin{proof}
  This is the declared model definition of displayed Process Kind.
\end{proof}
\begin{definition}[displayed Segment Orientation]
  \label{def:physics:phyx-mini-0425:phyxminiproblems-problemphyxmini0425-displayedsegmentorientation}
  \lean{PhyXMiniProblems.ProblemPhyXMini0425.displayedSegmentOrientation}
  \uses{def:physics:phyx-mini-0425:phyxminiproblems-problemphyxmini0425-cycleleg, def:physics:phyx-mini-0425:phyxminiproblems-problemphyxmini0425-segmentorientation}
  Segment orientations read from the primary image
\end{definition}
\begin{proof}
  This is the declared model definition of displayed Segment Orientation.
\end{proof}
\begin{definition}[Axis Quantity]
  \label{def:physics:phyx-mini-0425:phyxminiproblems-problemphyxmini0425-axisquantity}
  \lean{PhyXMiniProblems.ProblemPhyXMini0425.AxisQuantity}
  Physical quantity assigned to one of the two plotted axes
\end{definition}
\begin{proof}
  This is the declared model definition of Axis Quantity.
\end{proof}
\begin{definition}[Pressure Display Unit]
  \label{def:physics:phyx-mini-0425:phyxminiproblems-problemphyxmini0425-pressuredisplayunit}
  \lean{PhyXMiniProblems.ProblemPhyXMini0425.PressureDisplayUnit}
  Pressure unit printed beside the vertical axis
\end{definition}
\begin{proof}
  This is the declared model definition of Pressure Display Unit.
\end{proof}
\begin{definition}[Diagram Label]
  \label{def:physics:phyx-mini-0425:phyxminiproblems-problemphyxmini0425-diagramlabel}
  \lean{PhyXMiniProblems.ProblemPhyXMini0425.DiagramLabel}
  Literal labels and numeric annotations visible in the primary image
\end{definition}
\begin{proof}
  This is the declared model definition of Diagram Label.
\end{proof}
\begin{definition}[Rectangular PVCycle Figure]
  \label{def:physics:phyx-mini-0425:phyxminiproblems-problemphyxmini0425-rectangularpvcyclefigure}
  \lean{PhyXMiniProblems.ProblemPhyXMini0425.RectangularPVCycleFigure}
  \uses{def:physics:phyx-mini-0425:phyxminiproblems-problemphyxmini0425-volumequantity, def:physics:phyx-mini-0425:phyxminiproblems-problemphyxmini0425-cyclestate, def:physics:phyx-mini-0425:phyxminiproblems-problemphyxmini0425-cycleleg, def:physics:phyx-mini-0425:phyxminiproblems-problemphyxmini0425-axisquantity, def:physics:phyx-mini-0425:phyxminiproblems-problemphyxmini0425-pressuredisplayunit, def:physics:phyx-mini-0425:phyxminiproblems-problemphyxmini0425-diagramlabel}
  The rectangular pressure--volume diagram and its physical coordinates
\end{definition}
\begin{proof}
  This is the declared model definition of Rectangular PVCycle Figure.
\end{proof}
\begin{definition}[Gas Model]
  \label{def:physics:phyx-mini-0425:phyxminiproblems-problemphyxmini0425-gasmodel}
  \lean{PhyXMiniProblems.ProblemPhyXMini0425.GasModel}
  Material model specified in the problem
\end{definition}
\begin{proof}
  This is the declared model definition of Gas Model.
\end{proof}
\begin{definition}[Thermodynamic Device Role]
  \label{def:physics:phyx-mini-0425:phyxminiproblems-problemphyxmini0425-thermodynamicdevicerole}
  \lean{PhyXMiniProblems.ProblemPhyXMini0425.ThermodynamicDeviceRole}
  Role of the cyclic thermodynamic device
\end{definition}
\begin{proof}
  This is the declared model definition of Thermodynamic Device Role.
\end{proof}
\begin{definition}[Monatomic Rectangular Heat Engine]
  \label{def:physics:phyx-mini-0425:phyxminiproblems-problemphyxmini0425-monatomicrectangularheatengine}
  \lean{PhyXMiniProblems.ProblemPhyXMini0425.MonatomicRectangularHeatEngine}
  \uses{def:physics:phyx-mini-0425:phyxminiproblems-problemphyxmini0425-cyclestate, def:physics:phyx-mini-0425:phyxminiproblems-problemphyxmini0425-cycleleg, def:physics:phyx-mini-0425:phyxminiproblems-problemphyxmini0425-processkind, def:physics:phyx-mini-0425:phyxminiproblems-problemphyxmini0425-segmentorientation, def:physics:phyx-mini-0425:phyxminiproblems-problemphyxmini0425-rectangularpvcyclefigure, def:physics:phyx-mini-0425:phyxminiproblems-problemphyxmini0425-gasmodel, def:physics:phyx-mini-0425:phyxminiproblems-problemphyxmini0425-thermodynamicdevicerole}
  The working gas, its diagram, and its state and process observables. Amount of substance is kept abstract; only the explicitly named mole readout is a real number because Physlib's unit dimensions do not include amount of substance. Heat is positive into the gas and work is positive when done by the gas
\end{definition}
\begin{proof}
  This is the declared model definition of Monatomic Rectangular Heat Engine.
\end{proof}
\begin{definition}[amount Of Gas In Moles]
  \label{def:physics:phyx-mini-0425:phyxminiproblems-problemphyxmini0425-amountofgasinmoles}
  \lean{PhyXMiniProblems.ProblemPhyXMini0425.amountOfGasInMoles}
  \uses{def:physics:phyx-mini-0425:phyxminiproblems-problemphyxmini0425-monatomicrectangularheatengine}
  Mole readout of the physical gas sample
\end{definition}
\begin{proof}
  This is the declared model definition of amount Of Gas In Moles.
\end{proof}
\begin{definition}[gas Temperature In Kelvin]
  \label{def:physics:phyx-mini-0425:phyxminiproblems-problemphyxmini0425-gastemperatureinkelvin}
  \lean{PhyXMiniProblems.ProblemPhyXMini0425.gasTemperatureInKelvin}
  \uses{def:physics:phyx-mini-0425:phyxminiproblems-problemphyxmini0425-temperatureinkelvin, def:physics:phyx-mini-0425:phyxminiproblems-problemphyxmini0425-cyclestate, def:physics:phyx-mini-0425:phyxminiproblems-problemphyxmini0425-monatomicrectangularheatengine}
  Kelvin readout at a numbered state
\end{definition}
\begin{proof}
  This is the declared model definition of gas Temperature In Kelvin.
\end{proof}
\begin{definition}[internal Energy In Joules]
  \label{def:physics:phyx-mini-0425:phyxminiproblems-problemphyxmini0425-internalenergyinjoules}
  \lean{PhyXMiniProblems.ProblemPhyXMini0425.internalEnergyInJoules}
  \uses{def:physics:phyx-mini-0425:phyxminiproblems-problemphyxmini0425-energyinjoules, def:physics:phyx-mini-0425:phyxminiproblems-problemphyxmini0425-cyclestate, def:physics:phyx-mini-0425:phyxminiproblems-problemphyxmini0425-monatomicrectangularheatengine}
  Joule readout of the internal energy at a numbered state
\end{definition}
\begin{proof}
  This is the declared model definition of internal Energy In Joules.
\end{proof}
\begin{definition}[work Done By Gas In Joules]
  \label{def:physics:phyx-mini-0425:phyxminiproblems-problemphyxmini0425-workdonebygasinjoules}
  \lean{PhyXMiniProblems.ProblemPhyXMini0425.workDoneByGasInJoules}
  \uses{def:physics:phyx-mini-0425:phyxminiproblems-problemphyxmini0425-energyinjoules, def:physics:phyx-mini-0425:phyxminiproblems-problemphyxmini0425-cycleleg, def:physics:phyx-mini-0425:phyxminiproblems-problemphyxmini0425-monatomicrectangularheatengine}
  Signed joule readout of work done by the gas on a directed leg
\end{definition}
\begin{proof}
  This is the declared model definition of work Done By Gas In Joules.
\end{proof}
\begin{definition}[heat Transferred Into Gas In Joules]
  \label{def:physics:phyx-mini-0425:phyxminiproblems-problemphyxmini0425-heattransferredintogasinjoules}
  \lean{PhyXMiniProblems.ProblemPhyXMini0425.heatTransferredIntoGasInJoules}
  \uses{def:physics:phyx-mini-0425:phyxminiproblems-problemphyxmini0425-energyinjoules, def:physics:phyx-mini-0425:phyxminiproblems-problemphyxmini0425-cycleleg, def:physics:phyx-mini-0425:phyxminiproblems-problemphyxmini0425-monatomicrectangularheatengine}
  Signed joule readout of heat transferred into the gas on a leg
\end{definition}
\begin{proof}
  This is the declared model definition of heat Transferred Into Gas In Joules.
\end{proof}
\begin{definition}[Matches Problem Statement]
  \label{def:physics:phyx-mini-0425:phyxminiproblems-problemphyxmini0425-matchesproblemstatement}
  \lean{PhyXMiniProblems.ProblemPhyXMini0425.MatchesProblemStatement}
  \uses{def:physics:phyx-mini-0425:phyxminiproblems-problemphyxmini0425-monatomicrectangularheatengine, def:physics:phyx-mini-0425:phyxminiproblems-problemphyxmini0425-amountofgasinmoles, def:physics:phyx-mini-0425:phyxminiproblems-problemphyxmini0425-heattransferredintogasinjoules}
  The data stated in the prose. This records the measured heat supplied on 1 → 2, but neither net work, total heat input, nor thermal efficiency
\end{definition}
\begin{proof}
  This is the declared model definition of Matches Problem Statement.
\end{proof}
\begin{definition}[Matches Primary PVDiagram]
  \label{def:physics:phyx-mini-0425:phyxminiproblems-problemphyxmini0425-matchesprimarypvdiagram}
  \lean{PhyXMiniProblems.ProblemPhyXMini0425.MatchesPrimaryPVDiagram}
  \uses{def:physics:phyx-mini-0425:phyxminiproblems-problemphyxmini0425-volumeincubicmeters, def:physics:phyx-mini-0425:phyxminiproblems-problemphyxmini0425-pressureinkilopascals, def:physics:phyx-mini-0425:phyxminiproblems-problemphyxmini0425-displayedprocesskind, def:physics:phyx-mini-0425:phyxminiproblems-problemphyxmini0425-displayedsegmentorientation, def:physics:phyx-mini-0425:phyxminiproblems-problemphyxmini0425-monatomicrectangularheatengine, def:physics:phyx-mini-0425:phyxminiproblems-problemphyxmini0425-gastemperatureinkelvin}
  Exact transcription of the primary bitmap: clockwise arrows around a rectangle, vertical isochoric legs, horizontal isobaric legs, lower pressure 300 kPa, T₁ = 300 K, and Vmax = 2 Vmin. The value of pmax is not given here and must be derived from the heating process
\end{definition}
\begin{proof}
  This is the declared model definition of Matches Primary PVDiagram.
\end{proof}
\begin{definition}[Uses Classroom SIGas Constant]
  \label{def:physics:phyx-mini-0425:phyxminiproblems-problemphyxmini0425-usesclassroomsigasconstant}
  \lean{PhyXMiniProblems.ProblemPhyXMini0425.UsesClassroomSIGasConstant}
  \uses{def:physics:phyx-mini-0425:phyxminiproblems-problemphyxmini0425-monatomicrectangularheatengine}
  The standard classroom SI calibration R = 8.31 J/(mol K) used at the precision of the problem data and answer choices. It is an independent physical-constant readout, not an efficiency assumption
\end{definition}
\begin{proof}
  This is the declared model definition of Uses Classroom SIGas Constant.
\end{proof}
\begin{definition}[Has Physical Cycle Parameters]
  \label{def:physics:phyx-mini-0425:phyxminiproblems-problemphyxmini0425-hasphysicalcycleparameters}
  \lean{PhyXMiniProblems.ProblemPhyXMini0425.HasPhysicalCycleParameters}
  \uses{def:physics:phyx-mini-0425:phyxminiproblems-problemphyxmini0425-volumeincubicmeters, def:physics:phyx-mini-0425:phyxminiproblems-problemphyxmini0425-pressureinpascals, def:physics:phyx-mini-0425:phyxminiproblems-problemphyxmini0425-monatomicrectangularheatengine, def:physics:phyx-mini-0425:phyxminiproblems-problemphyxmini0425-amountofgasinmoles, def:physics:phyx-mini-0425:phyxminiproblems-problemphyxmini0425-gastemperatureinkelvin}
  Positivity and nondegeneracy conditions for a physical engine cycle
\end{definition}
\begin{proof}
  This is the declared model definition of Has Physical Cycle Parameters.
\end{proof}
\begin{definition}[Satisfies Monatomic Ideal Gas Laws]
  \label{def:physics:phyx-mini-0425:phyxminiproblems-problemphyxmini0425-satisfiesmonatomicidealgaslaws}
  \lean{PhyXMiniProblems.ProblemPhyXMini0425.SatisfiesMonatomicIdealGasLaws}
  \uses{def:physics:phyx-mini-0425:phyxminiproblems-problemphyxmini0425-volumeincubicmeters, def:physics:phyx-mini-0425:phyxminiproblems-problemphyxmini0425-pressureinpascals, def:physics:phyx-mini-0425:phyxminiproblems-problemphyxmini0425-legsource, def:physics:phyx-mini-0425:phyxminiproblems-problemphyxmini0425-legtarget, def:physics:phyx-mini-0425:phyxminiproblems-problemphyxmini0425-monatomicrectangularheatengine, def:physics:phyx-mini-0425:phyxminiproblems-problemphyxmini0425-amountofgasinmoles, def:physics:phyx-mini-0425:phyxminiproblems-problemphyxmini0425-gastemperatureinkelvin, def:physics:phyx-mini-0425:phyxminiproblems-problemphyxmini0425-internalenergyinjoules, def:physics:phyx-mini-0425:phyxminiproblems-problemphyxmini0425-workdonebygasinjoules, def:physics:phyx-mini-0425:phyxminiproblems-problemphyxmini0425-heattransferredintogasinjoules}
  Macroscopic governing laws for the monatomic ideal-gas cycle: p V = n R T at each equilibrium state; U = (3/2) n R T for a monatomic ideal gas; zero boundary work on an isochoric leg; W\_by = p (V\_f - V\_i) on an isobaric leg; Q\_in = U\_f - U\_i + W\_by on every leg. These laws are general relations. No field fixes a leg-energy result other than the measured 1 → 2 heat, net cycle work, total heat input, efficiency, or answer choice
\end{definition}
\begin{proof}
  This is the declared model definition of Satisfies Monatomic Ideal Gas Laws.
\end{proof}
\begin{definition}[net Cycle Work In Joules]
  \label{def:physics:phyx-mini-0425:phyxminiproblems-problemphyxmini0425-netcycleworkinjoules}
  \lean{PhyXMiniProblems.ProblemPhyXMini0425.netCycleWorkInJoules}
  \uses{def:physics:phyx-mini-0425:phyxminiproblems-problemphyxmini0425-cycleleg, def:physics:phyx-mini-0425:phyxminiproblems-problemphyxmini0425-monatomicrectangularheatengine, def:physics:phyx-mini-0425:phyxminiproblems-problemphyxmini0425-workdonebygasinjoules}
  Net work done by the gas over one complete traversal, in joules
\end{definition}
\begin{proof}
  This is the declared model definition of net Cycle Work In Joules.
\end{proof}
\begin{definition}[total Heat Input In Joules]
  \label{def:physics:phyx-mini-0425:phyxminiproblems-problemphyxmini0425-totalheatinputinjoules}
  \lean{PhyXMiniProblems.ProblemPhyXMini0425.totalHeatInputInJoules}
  \uses{def:physics:phyx-mini-0425:phyxminiproblems-problemphyxmini0425-cycleleg, def:physics:phyx-mini-0425:phyxminiproblems-problemphyxmini0425-monatomicrectangularheatengine, def:physics:phyx-mini-0425:phyxminiproblems-problemphyxmini0425-heattransferredintogasinjoules}
  Total heat entering the gas, obtained by summing positive leg heats
\end{definition}
\begin{proof}
  This is the declared model definition of total Heat Input In Joules.
\end{proof}
\begin{definition}[thermal Efficiency]
  \label{def:physics:phyx-mini-0425:phyxminiproblems-problemphyxmini0425-thermalefficiency}
  \lean{PhyXMiniProblems.ProblemPhyXMini0425.thermalEfficiency}
  \uses{def:physics:phyx-mini-0425:phyxminiproblems-problemphyxmini0425-monatomicrectangularheatengine, def:physics:phyx-mini-0425:phyxminiproblems-problemphyxmini0425-netcycleworkinjoules, def:physics:phyx-mini-0425:phyxminiproblems-problemphyxmini0425-totalheatinputinjoules}
  Dimensionless thermal efficiency W\_net / Q\_in
\end{definition}
\begin{proof}
  This is the declared model definition of thermal Efficiency.
\end{proof}
\begin{definition}[Answer Choice]
  \label{def:physics:phyx-mini-0425:phyxminiproblems-problemphyxmini0425-answerchoice}
  \lean{PhyXMiniProblems.ProblemPhyXMini0425.AnswerChoice}
  Labels of the four printed answer choices
\end{definition}
\begin{proof}
  This is the declared model definition of Answer Choice.
\end{proof}
\begin{definition}[displayed Efficiency]
  \label{def:physics:phyx-mini-0425:phyxminiproblems-problemphyxmini0425-displayedefficiency}
  \lean{PhyXMiniProblems.ProblemPhyXMini0425.displayedEfficiency}
  \uses{def:physics:phyx-mini-0425:phyxminiproblems-problemphyxmini0425-answerchoice}
  Fractional efficiency printed beside each answer label
\end{definition}
\begin{proof}
  This is the declared model definition of displayed Efficiency.
\end{proof}
\begin{definition}[recorded Answer Choice]
  \label{def:physics:phyx-mini-0425:phyxminiproblems-problemphyxmini0425-recordedanswerchoice}
  \lean{PhyXMiniProblems.ProblemPhyXMini0425.recordedAnswerChoice}
  \uses{def:physics:phyx-mini-0425:phyxminiproblems-problemphyxmini0425-answerchoice}
  Answer choice recorded in the supplied dataset metadata
\end{definition}
\begin{proof}
  This is the declared model definition of recorded Answer Choice.
\end{proof}
\begin{definition}[Rounds To Displayed Hundredth]
  \label{def:physics:phyx-mini-0425:phyxminiproblems-problemphyxmini0425-roundstodisplayedhundredth}
  \lean{PhyXMiniProblems.ProblemPhyXMini0425.RoundsToDisplayedHundredth}
  Compatibility with an efficiency displayed to the nearest hundredth
\end{definition}
\begin{proof}
  This is the declared model definition of Rounds To Displayed Hundredth.
\end{proof}
\begin{lemma}[cycle Leg Energy Readouts]
  \label{lem:physics:phyx-mini-0425:phyxminiproblems-problemphyxmini0425-cyclelegenergyreadouts}
  \lean{PhyXMiniProblems.ProblemPhyXMini0425.cycleLegEnergyReadouts}
  \uses{def:physics:phyx-mini-0425:phyxminiproblems-problemphyxmini0425-monatomicrectangularheatengine, def:physics:phyx-mini-0425:phyxminiproblems-problemphyxmini0425-workdonebygasinjoules, def:physics:phyx-mini-0425:phyxminiproblems-problemphyxmini0425-heattransferredintogasinjoules, def:physics:phyx-mini-0425:phyxminiproblems-problemphyxmini0425-matchesproblemstatement, def:physics:phyx-mini-0425:phyxminiproblems-problemphyxmini0425-matchesprimarypvdiagram, def:physics:phyx-mini-0425:phyxminiproblems-problemphyxmini0425-usesclassroomsigasconstant, def:physics:phyx-mini-0425:phyxminiproblems-problemphyxmini0425-hasphysicalcycleparameters, def:physics:phyx-mini-0425:phyxminiproblems-problemphyxmini0425-satisfiesmonatomicidealgaslaws}
  The governing laws and source data determine all signed leg energies. This is an intermediate derived result; none of these values is a premise field
\end{lemma}
\begin{proof}
  The conclusion follows from the governing relations and figure readouts named in the statement of cycle Leg Energy Readouts.
\end{proof}
\begin{lemma}[net Work And Heat Input Readouts]
  \label{lem:physics:phyx-mini-0425:phyxminiproblems-problemphyxmini0425-networkandheatinputreadouts}
  \lean{PhyXMiniProblems.ProblemPhyXMini0425.netWorkAndHeatInputReadouts}
  \uses{def:physics:phyx-mini-0425:phyxminiproblems-problemphyxmini0425-monatomicrectangularheatengine, def:physics:phyx-mini-0425:phyxminiproblems-problemphyxmini0425-matchesproblemstatement, def:physics:phyx-mini-0425:phyxminiproblems-problemphyxmini0425-matchesprimarypvdiagram, def:physics:phyx-mini-0425:phyxminiproblems-problemphyxmini0425-usesclassroomsigasconstant, def:physics:phyx-mini-0425:phyxminiproblems-problemphyxmini0425-hasphysicalcycleparameters, def:physics:phyx-mini-0425:phyxminiproblems-problemphyxmini0425-satisfiesmonatomicidealgaslaws, def:physics:phyx-mini-0425:phyxminiproblems-problemphyxmini0425-netcycleworkinjoules, def:physics:phyx-mini-0425:phyxminiproblems-problemphyxmini0425-totalheatinputinjoules}
  Consequently the engine does 2500 J of net work and absorbs 32465/2 J of heat per cycle. These are conclusions used by the final efficiency theorem
\end{lemma}
\begin{proof}
  The conclusion follows from the governing relations and figure readouts named in the statement of net Work And Heat Input Readouts.
\end{proof}
% --- Archon declaration topology end ---
% --- Archon physics formalization source end ---
